\documentclass[11pt]{article}
\usepackage[T1]{fontenc}
\usepackage{lmodern,microtype,amsmath,amssymb,amsthm,mathtools,booktabs}
\usepackage[letterpaper,margin=0.95in]{geometry}
\usepackage{xcolor}
\definecolor{ink}{HTML}{164F5A}
\usepackage[colorlinks=true,linkcolor=ink,citecolor=ink,urlcolor=ink,pdfauthor={OpenAI},pdftitle={Two elliptic traces determine the four-step tower distribution}]{hyperref}
\usepackage{fancyhdr}
\pagestyle{fancy}\fancyhf{}\fancyhead[L]{\small\itshape Four-step Cayley towers}\fancyhead[R]{\small Arithmetic of the moment curve}\fancyfoot[C]{\thepage}\setlength{\headheight}{14pt}
\newtheorem{theorem}{Theorem}[section]
\newtheorem{corollary}[theorem]{Corollary}
\newtheorem{lemma}[theorem]{Lemma}
\newcommand{\F}{\mathbb F}\newcommand{\Q}{\mathbb Q}\newcommand{\PP}{\mathbb P}
\newcommand{\Z}{\mathbb Z}\newcommand{\Jac}{\operatorname{Jac}}
\newcommand{\sq}{\operatorname{Pairs}}
\newcommand{\gh}{https://github.com/VladimirReshetnikov/ProveIt/blob/6d284fcf75bc3a91bb2e79ba7636f94724bff62c/Analysis/GeneralizedRiemannHypothesis}
\setlength{\parskip}{2pt}
\title{\vspace*{-0.45in}Two elliptic traces determine the distribution of spanning-tree valuations in four-step towers}
\author{Research note for the ProveIt project}
\date{27 November 2026\vspace{-0.13in}}
\begin{document}
\maketitle
\begin{abstract}
For every prime $p>8$ we compute the full frequency distribution of the
$p$-adic spanning-tree valuation of a rank-two abelian Cayley tower with
four ordered steps, simultaneously at every level. The arithmetic inputs
are the Frobenius traces of two fixed elliptic curves over $\Q$. We give
explicit quotient maps from the genus-nine moment curve, identify both
rational elliptic factors and their diagonal twists, and count all row
planes by secants, tangents and lines on a split quadric. Among the
$2(p+1)$ totally isotropic planes, the count in the middle valuation
class fluctuates around $(p+1)/2$ on the scale $\sqrt p$. Its exact
limiting law is a weighted sum of two independent semicircle variables.
\end{abstract}

\section{From tower valuations to an arithmetic curve}
A spanning list of four column vectors $s_i\in\Z^2$ means that their
reductions span $\F_p^2$. On $(\Z/p^n\Z)^2$, insert the undirected
edge $\{v,v+s_i\}$ for each pair $(v,i)$, retaining multiplicities.
Write $\kappa_n$ for the number of spanning trees. We count the
two-dimensional row spaces $W\subset\F_p^4$ of the step matrices.
Every such space has exactly $|\mathrm{GL}_2(\F_p)|$ ordered row bases,
so this census also gives the probability law for a uniformly random
spanning residue matrix. Lists can have zero or repeated columns.
The count is before quotienting by column permutations or signs.

For $P=[x_1:x_2:x_3:x_4]$ define the directional depth
$r(P)=\min\{j\ge1:\sum_i x_i^{2j}\ne0\}$. It is at most four by
Newton's identities applied to the four squares. The repository's
moment formula for $p>8$ is
\begin{equation}\label{eq:tower}
 A(W)=2\sum_{P\in\PP(W)}r(P),\qquad
 v_p(\kappa_n)= A(W)\frac{p^n-1}{p-1}-2n .
\end{equation}
It results from comparing the first nonzero terms of the Laplacian
eigenvalues at prime-power roots of unity \cite{article}. Thus a
single census of $A$ determines every level. Set
\[
 Q:\ \sum_i x_i^2=0,\qquad C:\ \sum_i x_i^2=\sum_i x_i^4=0
 \quad\text{in }\PP^3.
\]
Our first result evaluates the curve count in the repository's
four-step calculation. The second (Theorem~\ref{th:all}) computes
the census, including the row spaces on which the quadratic moment
is nonzero.

\begin{theorem}\label{th:curves}
Put $E_1:y^2=x(x+2)(x-1)$ and $E_2:y^2=x(x-1)(x-4)$. Then
\[
 \Jac(C)\ \sim_{\Q}\ E_1^6\times E_2^3 .
\]
For every $p\ge5$, with $a_p=p+1-\#E_1(\F_p)$ and
$b_p=p+1-\#E_2(\F_p)$, the zeta function of $C/\F_p$ has numerator
\[
 (1-a_pT+pT^2)^6(1-b_pT+pT^2)^3 .
\]
In particular $c_p:=\#C(\F_p)=p+1-6a_p-3b_p$.
\end{theorem}

\paragraph{Relation to earlier work.}
ProveIt's merged article computes the four-step maximum and identifies
the genus-nine curve that controls it \cite{article}. Its subsequent
four-step report computes the three frequencies on totally isotropic
planes in terms of $c_p$ \cite{short}. Paulhus had already obtained
the geometric multiplicity pattern $3+6$ for a genus-nine curve
with automorphism group of order $192$, as well as the octahedral
genus-three elliptic cube \cite[\S\S3.2.4, 4.2]{paulhus}.
Arora et al. analyzed rational twists of that genus-three curve
\cite{arora}. The specific rational quotient maps, the resulting
tower frequencies for all row planes, and their fluctuations are
the contributions developed here. The argument is unconditional.
It uses the article's cyclotomic valuation method through
\eqref{eq:tower}.

\section{Explicit maps to two elliptic curves}
The following proof keeps track of the ground field, which is essential
for point counts. A geometric decomposition alone leaves the signs of
rational twists undetermined. The complete intersection $C$ is smooth
in characteristic zero and at $p\ge5$. Indeed, dependent gradients
would force all nonzero $x_i^2$ to equal some $c\ne0$ and $mc=0$ for
$1\le m\le4$. Over an algebraic closure it is a connected divisor of
bidegree $(4,4)$ on a smooth quadric (two disjoint effective parts
cannot sum to $(4,4)$ by the intersection pairing). Adjunction gives
genus nine.

Put $y_i=x_i^2$. The image is the conic
$B:\sum y_i=\sum y_i^2=0$. The cover has group
$G=\{\pm1\}^4/\{\pm(1,1,1,1)\}$ of order eight, as is seen on the
open set with no coordinate zero. Its seven nontrivial characters
are indexed by the six pairs and the full set of four indices. For
an even subset $S$ the corresponding quotient double cover has
equation $t^2=\prod_{i\in S}y_i$ over $B$, interpreted homogeneously
with $t$ of weight $|S|/2$. Each coordinate line meets $B$ in two
simple points, and the four pairs are disjoint. Indeed, with one
coordinate zero, the other three are proportional to the three cube
roots of unity in some order. Two zero coordinates are impossible.
Riemann--Hurwitz gives genus one for each pair cover and genus three
for the full cover $H$.

Here and below, quotient differentials pull back to the subspace
invariant under the quotient subgroup. This follows by descent for
regular differentials on smooth curves in characteristic zero. Since
the invariant part for $G$ itself is zero, the pair and full-cover
subspaces occupy distinct character spaces. Their dimensions add to
nine. The sum of the seven pullbacks on Jacobians is therefore an
isogeny, defined over $\Q$. Permutations over $\Q$ identify all six
pair covers.

Use the chart $y_1+y_2=1$ and let $v=y_1-y_2$, $w=y_3-y_4$.
Then $y_3+y_4=-1$ and $v^2+w^2=-2$. On the pair cover for
$\{1,2\}$ take $h=2t$ so $h^2=1-v^2$. The map
\[
 (v,w,h)\longmapsto (v^2,vwh)
\]
takes values on $E_1$ and is nonconstant. A rational map between
smooth projective curves extends to a morphism. The induced map
on genus-one Jacobians is a rational isogeny, even when the source
curve has no rational base point.

On $H$ take $h=4t$ so $h^2=(1-v^2)(v^2+3)$. The double
transposition $\sigma=(12)(34)$ negates both $v,w$ and fixes $h$.
With $s=v^2+1$ we get
\[
 (vwh)^2=(s^2-1)(s^2-4),\qquad
 (v,w,h)\longmapsto(s^2,s vwh)\in E_2 .
\]
This nonconstant map factors through $H/\langle\sigma\rangle$.
For completeness, $\sigma$ has exactly four geometric fixed points
on $H$. On the conic with homogeneous coordinates
$(u,v,w)=(y_1+y_2,y_1-y_2,y_3-y_4)$ it fixes exactly the two
points with $u=0$. Both lie off the branch divisor, and the lift
fixes both points above each because $t$ has even weight.
Riemann--Hurwitz gives genus one for the quotient. The three
nonidentity elements of the double-transposition Klein four group
$V$ are conjugate by rational permutations. Their quotient genera
are all one. If $m_0$ is the dimension of the $V$-invariant
differentials, and $m_i$ the three other character dimensions, then
$m_0+m_i=1$ and $m_0+\sum_i m_i=3$. Thus $m_0=0$ and $m_i=1$.
The three pullbacks give $\Jac(H)\sim_{\Q} E_2^3$, proving the
rational decomposition. Good reduction of all three curves and
invariance of the Frobenius polynomial under isogeny give the
remaining assertions of Theorem~\ref{th:curves}.

\paragraph{Which rational form occurs?}
On writing $(y_i)=(l+m+n,l-m-n,-l+m-n,-l-m+n)$, the conic
becomes $l^2+m^2+n^2=0$. On it $\prod y_i=2(l^4+m^4+n^4)$.
This identifies $H$ as the central quadratic twist by $-1$ of
the model in \cite[Eq.~(1.4.1)]{arora}, after rescaling the
coordinate of weight two. The explicit map above fixes the twist
of $E_2$ in the trace formula.

\begin{corollary}[All diagonal twists]\label{cor:twist}
For nonzero $d_i\in\Q$, define $C_{\boldsymbol d}$ by
$\sum d_i x_i^2=\sum d_i^2 x_i^4=0$. Write $E^{(u)}$ for a
quadratic twist by $u$. Then
\[
 \Jac(C_{\boldsymbol d})\ \sim_{\Q} \prod_{i<j} E_1^{(d_i d_j)}
 \ \times\big(E_2^{(d_1d_2d_3d_4)}\big)^3 .
\]
\end{corollary}
\begin{proof}
Over the field generated by the square roots of the $d_i$, the
substitution $x_i\mapsto\sqrt{d_i}x_i$ identifies the curves.
Its descent action is a coordinate sign change. On the factor for
$S$ this acts by the product of the signs indexed by $S$, since the
deck involution acts by $-1$ on its differentials, and hence on its
Jacobian (the quotient is the conic). This is exactly twisting by
$\prod_{i\in S}d_i$. The same sign acts on all three elliptic factors
of $\Jac(H)$.
\end{proof}
For $p\ge5$ with all $d_i$ units, write $k$ for the number of
nonsquare $d_i$ modulo $p$. The Frobenius trace on the twisted
Jacobian is respectively $6a_p+3b_p$, $-3b_p$, or $-2a_p+3b_p$
as $\min(k,4-k)=0,1,2$. In the middle case the entire sixfold
contribution disappears from the trace.

\section{The complete tower census}
Set $s=p+1$, $I=\mathbf1_{p\equiv1\bmod8}$, $c=c_p$ and $d=48I$.
The quadric $Q$ is split over $\F_p$. For example choose $u^2+v^2=-1$,
possible since two sets of $(p+1)/2$ shifted square values intersect.
The vectors $(1,0,u,v)$ and $(0,1,-v,u)$ span a totally isotropic
plane. Consequently $Q$ has $s^2$ rational points in an $s$ by $s$
grid, with $2s$ generator lines, two through each point. Let
$D=\{P:r(P)=4\}$. The types $Q\setminus C$, $C\setminus D$, $D$
have weights $r-1=1,2,3$, respectively. Encode them by the polynomial
\[
 F(z)=(s^2-c)z+(c-d)z^2+d z^3,\qquad
 \sq(f)(z)=\frac{f(z)^2-f(z^2)}{2}.
\]
Here $\sq$ counts unordered pairs of distinct weighted points.
For the four types of generator line, let $f_j$ record its points,
$n_j$ its multiplicity and $e_j$ its sum of weights.
\begin{center}
\begin{tabular}{@{}llrr@{}}\toprule
Type $j$ & $f_j(z)$ & $n_j$ & $e_j$\\\midrule
$0$ & $sz$ & $2s-12I-c/2$ & $s$\\
$B$ & $(s-4)z+4z^2$ & $(c-d)/2$ & $s+4$\\
$D_2$ & $(s-2)z+2z^3$ & $24I$ & $s+4$\\
$D_4$ & $(s-4)z+4z^3$ & $12I$ & $s+8$\\\bottomrule
\end{tabular}
\end{center}

\begin{theorem}[Distribution for every row plane]\label{th:all}
For $p>8$ let $G_p(z)=\sum_{\dim W=2}z^{A(W)/2-s}$. Then
\begin{equation}\label{eq:census}
 G_p(z)=\frac{p^2(p-1)^2}{2}+(p-1)F(z)+\sq(F)(z)
       +\sum_j n_j\big(z^{e_j}-\sq(f_j)(z)\big).
\end{equation}
In particular $[z^m]G_p$ is the number of row planes with
$A(W)=2(s+m)$. The only curve counts required to evaluate the
formula are those of $E_1$ and $E_2$.
\end{theorem}
\begin{proof}
We first recover the generator table, as in \cite{short}. At a
depth-four point the four squared coordinates are the roots of
$T^4-\lambda$ for some $\lambda\ne0$. Such points exist precisely
when $8\mid p-1$. In that case, if $g$ has order eight, they
are signed permutations and projective rescalings of
$a=(1,g,g^2,g^3)$. Normalizing the first coordinate gives
$3!\,2^3=48$ points. The two generators through $[a]$ are spanned
by $a$ and, respectively,
\[
 b=(-g,1,-g^3,g^2),\qquad b'=(-g,1,g^3,-g^2).
\]
Indeed the vectors in each pair are mutually orthogonal and isotropic,
and $b\cdot b'=2(1+g^2)\ne0$. On the first plane the coordinates
have the form $(U,V,g^2U,g^2V)$ for independent $U,V$. The quartic
is $2(U^4+V^4)$ with four rational roots, and the sextic vanishes.
On $Xa+Yb'$ the quartic is $24g^2X^2Y^2$ and both its roots
have depth four. Each $D$ point therefore lies on exactly one
generator of each of the last two types. Incidence counting gives
$n_{D_4}=d/4$ and $n_{D_2}=d/2$.

On a generator avoiding $D$, the quartic has zero or four rational
roots. To see this, take independent columns of a row-basis matrix
as the identity and write $[I_2\mid M]$ with $MM^{\mathsf T}=-I_2$.
By a column sign change arrange
$M=\left(\begin{smallmatrix}u&v\\v&-u\end{smallmatrix}\right)$.
The matrices $M$ and
$J=\left(\begin{smallmatrix}0&-1\\1&0\end{smallmatrix}\right)$
permute the columns up to signs. Projectively they generate a Klein
four group, since they anticommute and $J^2=M^2=(JM)^2=-I_2$.
Every even moment is invariant under the transpose actions on
covectors. A quartic zero fixed by a nonidentity element would have
eigenvalue $\lambda$ with $\lambda^2=-1$. Its sextic must vanish
because $\lambda^6=-1$, giving a point of $D$. Thus every rational
quartic zero on such a line has an orbit of four. The quartic does
not vanish identically there since $C$ is geometrically irreducible
of genus nine. Each point of $C\setminus D$ is on two generators
avoiding $D$ and each contributing generator has four such points.
This gives $n_B$ and the rest of the table.

Each row plane corresponds to a projective line. Its exponent in
$G_p$ is the sum of $r(P)-1$ over its rational points. On a
non-generator this sum comes only from its intersection with $Q$.
Every pair of distinct $Q$-points determines one secant unless its
line is a generator. No pair belongs to two distinct generators.
This gives $\sq(F)-\sum_jn_j\sq(f_j)$. At any $Q$-point its tangent
plane contains $p+1$ lines through the point. Removing its two
generators leaves $p-1$ tangents, giving $(p-1)F$.
Each $Q$-point is joined to $p^2$ other $Q$-points outside its two
generators, so the number of secants is $s^2p^2/2$. Subtract these,
$s^2(p-1)$ tangents and $2s$ generators from the total number
$(p^2+1)(p^2+p+1)$ of projective lines. This leaves $p^2(p-1)^2/2$
external lines. Adding the generators themselves proves
\eqref{eq:census}.
\end{proof}

\paragraph{Two different generator distributions.}
At $p=17$, $a_p=-6$, $b_p=2$, giving $c=d=48$.
There are $36$ generators, of which $24$ have $A=80$ and $12$ have
$A=88$. Every generator at this prime passes through $D$. The
$A=88$ planes give $v_{17}(\kappa_n)=88(17^n-1)/16-2n$ at all levels.
For comparison, $p=13$ gives $c=32$ and $d=0$, so $16$ of its $28$
generators have $A=64$ and $12$ have $A=56$.
For a random spanning residue matrix as $p\to\infty$, the external
lines alone account for a proportion tending to $1/2$, all at
$A=2s$. Ordinary secants through two points of $Q\setminus C$
account for another proportion tending to $1/2$, all at $A=2s+4$.
Indeed the other secants involve at least one of $c=O(p)$ points,
so are $O(p^3)$ among a total of order $p^4$.

\subsection*{Signed permutation orbits and a congruence modulo 192}
The two elliptic traces satisfy a joint congruence much finer than
parity. Define $J_p=1$ when $-2$ is a quadratic residue and both
$-1+\sqrt{-2}$ and $-1-\sqrt{-2}$ are squares in $\F_p$; set
$J_p=0$ otherwise.
\begin{theorem}\label{th:192}
For $p\ge5$,
\[
 c_p\equiv 32\,\mathbf1_{p\equiv1\bmod3}
       +48\,\mathbf1_{p\equiv1\bmod8}+96J_p\pmod{192}.
\]
For every modulus $M$, even among primes $p\equiv1\pmod{24M}$,
the two possible values $80,176\pmod{192}$ each have relative
density $1/2$.
\end{theorem}
\begin{proof}
Act on $C(\F_p)$ by projective signed permutations, a group of
order $8\cdot24=192$. For $y_i=x_i^2$ the polynomial with these
four roots is $T^4-e_3T+e_4$, by the two moment equations and
Newton's identities. If $e_3e_4\ne0$ a rescaling that preserves
the multiset of $y_i$ must have both its cube and its fourth power
equal to one. Thus it is the identity scalar on the squares.
When all four squares are distinct and nonzero this leaves a
trivial projective stabilizer, so such orbits contribute multiples
of $192$.

There are three disjoint exceptions. When $e_4=0$, the squares
are projectively $0,1,\omega,\omega^2$ with $\omega$ a primitive
cube root. There are four choices for the zero, two orderings
after scaling a specified nonzero coordinate to one, and four
remaining sign choices, giving $32$ points if $3\mid p-1$.
When $e_3=0$ there are the $48I$ points of depth four computed
above (the same roots-of-unity count works at $5,7$).
In the remaining repeated-root case the squares scale to
$1,1,-1+\sqrt{-2},-1-\sqrt{-2}$. Indeed the sum and sum of squares
force the last two values to be roots of $T^2+2T+3$. They are
distinct and different from $0,1$ for $p\ge5$. There are six
choices for the equal pair, two orders for the remaining values
and eight sign choices after normalization. This gives $96J_p$.

For the density, let $A=\Q(\sqrt{-2})$ and
$\alpha=-1+\sqrt{-2}$, $\beta=-1-\sqrt{-2}$. The splitting field
$K=A(\sqrt\alpha,\sqrt\beta)$ of $T^4+2T^2+3$ has Galois group
the dihedral group of order eight. In fact $\alpha,\beta$ are
independent square classes in $A$: their norms are $3$, and
$\alpha\beta=3$ is also a nonsquare in $A$. Conjugation exchanges
the two square roots and the two independent sign changes.
The commutator changes both signs, so its fixed field is
$L=\Q(\sqrt{-2},\sqrt3)$. This is the largest abelian subfield and
lies in $\Q(\zeta_{24})$. For $p\ge5$, $J_p=1$ exactly when $p$
splits completely in $K$, by reduction of the four roots. With
$F=\Q(\zeta_{24M})$ one has $K\cap F=L$. In the compositum $KF$
the subgroup fixing $F$ has order two and is central (restriction
to $K$ identifies it with the central commutator). Chebotarev's
density theorem gives equal relative densities for the two elements
among primes splitting in $F$. The congruence gives the stated
residues.
\end{proof}

\section{Fluctuations over the primes}
The three generator valuation coefficients are $4s$, $4s+8$,
$4s+16$. Write $N_p$ for the number in the middle class. Its
value is $n_B+n_{D_2}=c_p/2$.

\begin{corollary}\label{cor:st}
As $p>8$ ranges over primes, the values
\[
 T_p=\frac{N_p-(p+1)/2}{\sqrt p}
\]
have limiting distribution $-3X-\tfrac32Y$, where $X,Y$ are independent
with density $\sqrt{4-x^2}/(2\pi)$ on $[-2,2]$. Their prime-average
limit is zero and their limiting variance is $45/4$. Furthermore
\[
 \liminf_{p\to\infty}T_p=-9,\qquad \limsup_{p\to\infty}T_p=9 .
\]
The proportion of isotropic row planes in this class tends to $1/4$.
It is above $1/4$ on a set of primes of density $1/2$ and below it
on a set of the same density.
\end{corollary}
\begin{proof}
The two elliptic $j$-invariants are $21952/9$ and $35152/9$.
They are nonintegral, so neither curve has complex multiplication.
At $p=7$ the traces are $-4$ and $0$, obtained by evaluating the
two cubics at the seven residues. The first reduction is ordinary,
the second supersingular, ruling out even a geometric isogeny
since isogenies specialize at good primes and preserve this
distinction. Thus the associated non-CM weight-two newforms are
twist-inequivalent. Joint Sato--Tate \cite[Thm.~1.2]{thorner}
gives independent semicircle limits for $a_p/\sqrt p,b_p/\sqrt p$.
Now $T_p=-(6a_p+3b_p)/(2\sqrt p)$ exactly. Each variable in the
limit has mean zero and variance one. Their product measure gives
positive mass arbitrarily near each corner of its square, and the
Hasse bounds give $|T_p|\le9$. Symmetry and continuity of the limit
give the two density statements.
\end{proof}

\subsection*{Mean and variance at every rank}
There is also a consequence for the article's higher-rank question.
For $1\le k\le4$ take a uniformly random $k$-dimensional row space
of four step columns. Define $A=2\sum_{P\in\PP(W)}r(P)$,
$u_j=(p^j-1)/(p-1)$ for $j\ge0$, and
\[
 \rho=\frac{u_k}{u_4},\qquad
 \tau=\frac{u_k u_{k-1}}{u_4 u_3},\qquad
 U=s^2+c+d,\qquad V=s^2+3c+5d.
\]
Then, with $s=p+1,c=c_p,d=48I$ as above,
\[
 \mathbb E[A]=2u_k+2\rho U,\qquad
 \mathrm{Var}(A)=4\big[(\rho-\tau)V+(\tau-\rho^2)U^2\big].
\]
Indeed each fixed point occurs with probability $\rho$, and each
fixed ordered pair of distinct points with probability $\tau$, by
double counting and transitivity. The sums of the point weights and
of their squares are respectively $U$ and $V$ by the three types
on $Q$. Expanding the first and second moments yields the formulas.
This computes both statistics in particular for rank three, where
the entire distribution remains to be determined.

\paragraph{Verification and further questions.}
The accompanying Python program compares the census with direct
row-space enumeration and the two elliptic character sums with
point counts from the defining equations. These checks sample the
proof and allow its sign conventions and exceptional classes to be
reproduced. A next question is to determine the rank-three
four-step distribution. Each row space then cuts out a plane
section of $C$; the incidence of several special points on the
same plane becomes relevant. A second question is to determine
which higher moment curves permit equally explicit rational
elliptic factors.

\begin{thebibliography}{9}
\bibitem{article} ProveIt contributors. \emph{Beyond the Riemann
hypothesis for Dirichlet $L$-functions}, section on
\href{\gh/Article/sections/gra.tex}{graph towers and spanning-tree
valuations}. Repository snapshot \texttt{6d284fcf75bc}, 2026.
\bibitem{short} ProveIt research packet.
\href{\gh/Research/short_tower_distributions_2026-11-27/short_tower_distributions.tex}{\emph{Three- and four-step graph towers. Distributions and extremizer orbits}},
27 November 2026. See the four-step distribution and orbit theorem.
\bibitem{paulhus} J. Paulhus. \href{https://doi.org/10.4064/aa132-3-3}{Decomposing Jacobians of curves with extra automorphisms}.
\emph{Acta Arith.} \textbf{132} (2008), 231--244.
\bibitem{arora} S. Arora, V. Cantoral-Farf\'an, A. Landesman,
D. Lombardo and J. S. Morrow.
\href{https://arxiv.org/abs/1608.06784}{The twisting Sato--Tate group of the curve $y^2=x^8-14x^4+1$}.
\bibitem{thorner} J. Thorner.
\href{https://arxiv.org/abs/2002.10450}{Effective forms of the Sato--Tate conjecture}.
\emph{Res. Math. Sci.} \textbf{8} (2021), article 4.
\end{thebibliography}
\end{document}
